\section{Datasets}
% Describe MUTAG, ENZYMES, IMDB-MULTI: type of graphs, 
% nodes, labels.

% Any preprocessing steps.
The experiments in this work are conducted on three widely used benchmark datasets from \textbf{TUDataset}~\cite{morris2020tudataset}: MUTAG, ENZYMES and IMDB-MULTI.
Although these datasets are relatively small, they are highly diverse in structure and 
they cover domains such as chemistry, biology and social networks.
They are standard benchmarks in the graph representation learning literature~\cite{narayanangraph2vec,tsitsulin2018netlsd,xu2019gin}
and are commonly used to evaluate graph embedding and classification methods.
Their popularity also enables fair comparison with existing methods and published results.

\subsection{MUTAG}
\textbf{MUTAG}~\cite{debnath1991mutag} is a molecular graph dataset consisting of chemical compounds, where each graph represents a molecule,
nodes correspond to atoms and edges represent chemical bonds.
The task is to classify molecules into two classes according to their mutagenic effect on a bacterium.
MUTAG contains a relatively small number of graphs with discrete node labels representing atom types.
Graphs in MUTAG are typically small and sparse, making the dataset suitable for fast experimentation.
Despite its size, it is structurally rich and has been widely adopted in the literature as a sanity check for graph learning methods.

\subsection{ENZYMES}
The \textbf{ENZYMES}~\cite{borgwardt2005graphkernels} dataset is derived from protein tertiary structures and consists of graphs representing enzymes.
Nodes correspond to secondary structure elements, while edges indicate spatial proximity or interactions between residues.
Each graph is labeled according to the enzyme's functional class, resulting in a multiclass classification task with six classes.
So compared to MUTAG, ENZYMES exhibits higher structural variability and a more challenging classification setting.
This benchmark is considered difficult for many graph embedding methods, often exposing limitations in generalization and representation capacity.

\subsection{IMDB-MULTI}
\textbf{IMDB-MULTI}~\cite{yanardag2015deepgraphkernels} is a social network dataset constructed from movie collaboration data.
Each graph represents the ego network of actors appearing in a movie, where nodes correspond to actors and edges represent co-appearance relationships.
Unlike molecular and biological datasets, IMDB-MULTI does not provide node labels or attributes,
which makes classification entirely dependent on graph topology.
This creates a true challenge for graph embedding methods, as they must rely solely on structural patterns such as connectivity, degree distributions and subgraph configurations.